\scanpage{32}\mspage{24}
% @ top, continuation | text | au lieu d'une base
au lieu d'une base, il faut en plus se donner un ens. de cardinal 2. (donc le groupe des autom est multiplié par $\unsure{\mathfrak S|\mathscr S}_2$, le nb d'objets est divisé par 2 \ldots).

% @ upper left, bracketed note | text | Étudier position relative
\margin{Étudier position relative d'une racine et d'un triangle}% slanted
% @ upper middle, divider | drawing | short slightly curved horizontal stroke
\figure{\begin{tikzpicture}\draw(0,0)..controls(.4,-.04) and (.72,.04)..(1.02,0)..controls(1.13,-.01) and (1.28,.01)..(1.4,.04);\end{tikzpicture}}
% @ upper half, triplets | text | Étudions les triples
Étudions les triples $(t,t',b)$, où $(t,t')$ sont deux triangles disjoints et $b$ une base orthogonale à $t$ \uline{et} $t'$ -- ils correspondent à un ens. $I$ de card 6, avec 2 partitions \inserted{$\Pi,\Pi'$} de types $(2,2,2)$, sans classe commune. On trouve que ceci correspond à une structure hexagonale sur $I$ \inserted{à co-triangle inscrit marqué} -- avec comme partitions celles marquées aux traits \rule{1.3em}{.7pt} \inserted{(pour $t$)} et \tikz\draw[decorate,decoration={zigzag,segment length=2pt,amplitude=.5pt}](0,0)--(.45,0); \inserted{(pour $t'$)}. \deleted{\unsure{Les autres paires}} Si on donne $\{t,t'\}$ au lieu de $(t,t')$, on n'a pas à choisir de cotriangle dans l'hexagone.
% @ middle left, hexagon | diagram | vertices 1–6, alternate sides and diagonals
\diagram{\begin{tikzcd}[ampersand replacement=\&,execute at end picture={
\coordinate (a) at (1.45,1.78);\coordinate (b) at (.57,1.88);\coordinate (c) at (-.04,1.15);
\coordinate (d) at (.43,.30);\coordinate (e) at (1.41,.33);\coordinate (f) at (2.00,1.12);
\draw[line width=2pt](a)--(b) (c)--(d) (e)--(f);
\draw[decorate,decoration={zigzag,segment length=2pt,amplitude=.5pt}](b)--(c) (d)--(e) (f)--(a);
\draw(a)--(c)--(e)--cycle (b)--(d)--(f)--cycle;
\draw[dashed](a)--(d) (b)--(e) (c)--(f);
\node[above right]at(a){$1$};\node[above left]at(b){$2$};\node[left]at(c){$3$};\node[below left]at(d){$4$};\node[below]at(e){$5$};\node[right]at(f){$6$};
}]\vphantom{A}\end{tikzcd}}
\margin{penser (plutôt) à l'hexagone dual -- dont les \uline{cotés} sont les él. de $I$ \ldots}% slanted

% @ lower half, completing triangle | text | On trouve que, si t″
On trouve que, si $t''$ est le triangle qui complète $\{t,t'\}$ en un bitriangle $\{t,t',t''\}$, alors $b$ est aussi orth. à $t''$ -- en fait on a la situation d'une base orthogonale à un bitriangle cf plus bas. Pour $b$ fixé, le syst. des $(t,t',t'')$ correspondants est de cardinal $120=5!$ (60 pour $\{t,t'\}$ $=\frac{5!}{2}=$ nb des structures hexagonales sur $I$ de card. 6) en fait en tout \deleted{$120.72=$} $72.5!=\operatorname{card}W/6$ \inserted{$\downarrow$} tels triples]
\margin{$(6=\operatorname{card}\mathfrak S_{\unsure{6|t}}$}% slanted
% @ foot, second count | text | ou nb
ou nb \inserted{(syst aussi :) }
% @ foot, three factors with captions | display | 45 . 32 . 6
\[
\underbrace{45}_{\text{triangles }t}\,.\,\underbrace{32}_{\text{triangles }t'\text{ disj. de }t}\,.\,
\underbrace{6}_{\substack{\text{nb de racines (ou bases)}\\\text{orth. à }t,t'}}
\]
% archive sheet number: 31
